\begin{table}[htbp]\centering\small
\caption{Сравнение SCC: наблюдаемые исходы.}
\label{tab:scc2000-quality}
\begin{tabular}{lrrrrr}\toprule
Метод & Назначено & Получено & Оценено & Успехов & Успех (\%)\\\midrule
SN & 2000 & 1998 & 1965 & 1064 & 54.15\\
SR & 2000 & 1997 & 1964 & 1067 & 54.33\\
SCC & 2000 & 1600 & 1574 & 809 & 51.40\\
\bottomrule\end{tabular}
\par\smallskip\noindent Оценено: пригодные исходы отдельных повторов. Процент равен числу успехов, делённому на число оценённых повторов; он не является парной оценкой с равными весами задач. Прерванные процессы остаются недоступными.
\end{table}

\begin{table}[htbp]\centering\small
\caption{Парные разности качества (процентные пункты).}
\label{tab:scc2000-contrasts}
\begin{tabular}{lrrrr}\toprule
Сравнение & Задачи & Разность & 95\% интервал & $p$ Холма\\\midrule
SCC -- SR & 883 & -2.68 & $[-4.72, -0.60]$ & 0.0217\\
SCC -- SN & 882 & -2.48 & $[-4.50, -0.43]$ & 0.0217\\
\bottomrule\end{tabular}
\par\smallskip\noindent Выбранные парные повторы сначала усредняются внутри каждой задачи; веса задач равны. Два теста образуют отдельное семейство Холма. Бутстрэп-интервалы 95\% являются описательными. Тире обозначает недоступную или неоцениваемую величину.
\end{table}

\begin{table}[htbp]\centering\small
\caption{Границы по назначенным исходам (процентные пункты).}
\label{tab:scc2000-bounds}
\begin{tabular}{lrrr}\toprule
Сравнение и выборка & Задачи & Границы & С пропусками\\\midrule
SCC -- SR (Пригодные) & 941 & $[-12.98, 6.61]$ & 340\\
SCC -- SR (Все выбранные) & 956 & $[-14.35, 8.07]$ & 355\\
SCC -- SN (Пригодные) & 941 & $[-13.14, 6.43]$ & 340\\
SCC -- SN (Все выбранные) & 956 & $[-14.50, 7.90]$ & 355\\
\bottomrule\end{tabular}
\par\smallskip\noindent В знаменателе каждой задачи сохраняется число выбранных назначенных повторов. Это границы возможных эффектов, а не доверительные интервалы.
\end{table}

\begin{table}[htbp]\centering\small
\caption{Отношение средних парных оценок стоимости.}
\label{tab:scc2000-resource-pooled}
\begin{tabular}{lrrr}\toprule
Сравнение & Задачи & Отношение & 95\% интервал\\\midrule
SCC / SR & 896 & 3.138 & $[3.07, 3.21]$\\
SCC / SN & 895 & 2.939 & $[2.87, 3.01]$\\
\bottomrule\end{tabular}
\par\smallskip\noindent В описательные оценки входят только парные завершённые процессы с полными счётчиками. Оценка по API-тарифам не равна платежу по подписке; известные частичные расходы учитываются отдельно.
\end{table}

\begin{table}[htbp]\centering\small
\caption{Среднее отношений стоимости по задачам.}
\label{tab:scc2000-resource-mean_ratio}
\begin{tabular}{lrrr}\toprule
Сравнение & Задачи & Отношение & 95\% интервал\\\midrule
SCC / SR & 896 & 3.424 & $[3.35, 3.49]$\\
SCC / SN & 895 & 3.234 & $[3.17, 3.30]$\\
\bottomrule\end{tabular}
\par\smallskip\noindent В описательные оценки входят только парные завершённые процессы с полными счётчиками. Оценка по API-тарифам не равна платежу по подписке; известные частичные расходы учитываются отдельно.
\end{table}

\begin{table}[htbp]\centering\small
\caption{Чувствительность к группам сходных задач (процентные пункты).}
\label{tab:scc2000-source}
\begin{tabular}{lrrr}\toprule
Группировка и сравнение & Группы & Задачи & 95\% интервал\\\midrule
Code + prompt 0.70; SCC -- SR & 876 & 883 & $[-4.72, -0.58]$\\
Code + prompt 0.70; SCC -- SN & 875 & 882 & $[-4.51, -0.43]$\\
Prompt 0.50; SCC -- SR & 880 & 883 & $[-4.76, -0.59]$\\
Prompt 0.50; SCC -- SN & 879 & 882 & $[-4.59, -0.42]$\\
Prompt 0.70; SCC -- SR & 883 & 883 & $[-4.72, -0.60]$\\
Prompt 0.70; SCC -- SN & 882 & 882 & $[-4.50, -0.43]$\\
Prompt 0.90; SCC -- SR & 883 & 883 & $[-4.72, -0.60]$\\
Prompt 0.90; SCC -- SN & 882 & 882 & $[-4.50, -0.43]$\\
\bottomrule\end{tabular}
\par\smallskip\noindent Целиком перевыбираются группы с сохранением равных весов задач. Эти интервалы не образуют дополнительных тестов гипотез.
\end{table}

\begin{table}[htbp]\centering\small
\caption{Диагностика этапов выполнения (процентные пункты).}
\label{tab:scc2000-timing}
\begin{tabular}{lrr}\toprule
Этап и сравнение & Задачи & Разность\\\midrule
До; SCC -- SR & 800 & -3.10\\
До; SCC -- SN & 800 & -2.48\\
После; SCC -- SR & 283 & -1.41\\
После; SCC -- SN & 282 & -2.42\\
На границе; SCC -- SR & 0 & ---\\
На границе; SCC -- SN & 0 & ---\\
\bottomrule\end{tabular}
\par\smallskip\noindent Этапы определены относительно исходной поправки о 2000 блоках, а не последующих возобновлений транспорта. Это описательные оценки; они не выделяют причинный эффект времени.
\end{table}
